\begin{figure}[htbp]
\centering
\begin{tikzpicture}[>=stealth,scale=2]
\draw[->,black!45] (-1.3,0)--(1.4,0) node[right] {$\operatorname{Re}z$};
\draw[->,black!45] (0,-1.3)--(0,1.4) node[above] {$\operatorname{Im}z$};
\draw[black!40] (0,0) circle(1);
\foreach \k in {0,1,2,3,4} {
 \draw[arkivblaa,thin] (0,0)--({cos(72*\k)},{sin(72*\k)});
 \fill[arkivblaa] ({cos(72*\k)},{sin(72*\k)}) circle(.014);
 \ifnum\k=0
  \node[font=\small,above] at (1.08,0) {$z_0$};
 \else
  \node[font=\small] at ({1.16*cos(72*\k)},{1.16*sin(72*\k)}) {$z_{\k}$};
 \fi
}
\draw[->,arkivblaa] (.28,0) arc[start angle=0,end angle=72,radius=.28];
\node at (.39,.28) {$\frac{2\pi}{5}$};
\end{tikzpicture}
\caption{De fem rødder i $z^5=1$: $z_k=e^{2\pi ik/5}$, $k=0,\ldots,4$. De ligger jævnt fordelt på enhedscirklen. For $w=re^{i\theta}\ne0$ skaleres radius til $r^{1/n}$ og vinklerne til $(\theta+2\pi k)/n$. Dimensionsløse akser med samme målestok.}
\label{fig:b1-komplekseroedder}
\end{figure}
